\section{The complete degree-three cochain}
\label{app:cochain}

This appendix specifies the finite cochain used in
Lemma~\ref{coc:data}, including the coefficient identities that verify
its properties.  Lower-case letters are the basis of $C$ from
\eqref{alg:C-corners}; capital letters are the dual basis of $DC$.
A four-letter entry $abcv$ in the table means that the indicated
coefficient is the coefficient of $v$ in $p(a,b,c)$.  All unlisted
coefficients are zero.  For example, the entry \texttt{tyJf} in the
$q^3$ row means $p(t,y,J)=q^3f$.
Each entry has composable radical inputs
and an output with the same two endpoints.

\begin{center}
\small
\begin{tabular}{r@{\quad}l}
 coefficient & input and output letters \\ \hline
 $q^2$ & \texttt{EuNT EyZX FNtU JtXZ NtEU NtXV XzJT ZxEY tEun tXuf} \\
 $q^2$ & \texttt{tXvn tXzj tuFn tuNf tvNn uFNT vNNT vtJu xEyz zJtx} \\
 $q$ & \texttt{EuUE EuVX EvVE ExXE EyYE EzZE FNnF FVvF FVxU FtTF} \\
 $q$ & \texttt{FtYU Ftun Ftxj JFtX JNtZ JUyX JVyZ JjEE JjTT JjXX} \\
 $q$ & \texttt{NFtU NNtV NUyU NVyV NjJN NjZV NnNN NnVV Ntuf Ntvn} \\
 $q$ & \texttt{Ntzj TtEE TtXX UuFF UuUU UvNF UvVU UxXU Uyun Uyxj} \\
 $q$ & \texttt{UzZU VvNN VvVV VxEU VxXV Vyuf Vyvn Vyzj VzZV XxEE} \\
 $q$ & \texttt{XxXX YuFT YuUY YvNT YvVY YxXY YyEE YyJJ YyTT YyXX} \\
 $q$ & \texttt{YyYY YyZZ YzZY ZxXZ ZzEE ZzJJ ZzTT ZzXX jJFF jJUU} \\
 $q$ & \texttt{jJjj jJnn jYyj jZvn jZzj nNFF nNUU ntTn tJjt tTNN} \\
 $q$ & \texttt{tTVV tTtt tYyt tZzt tzJn uFVX uFty uNFT uNUY uNte} \\
 $q$ & \texttt{uUyy uVye ujJu vNUE vNVX vNty vVyy vtTv xEuv xJFT} \\
 $q$ & \texttt{xJUY xJjx xJte xTtx xXuu xXvv xXyy xXzz xYyx xZye} \\
 $q$ & \texttt{xZzx xuFv xuNu xvNv yxJu yzJv zJNT zJVY zYyz zZuu} \\
 $q$ & \texttt{zZvv zZyy zZzz} \\
 $q^3$ & \texttt{EyJT JtEY tEyj tyJf} \\
 $1$ & \texttt{JNjX JVzX NNjU NVzU NjYU Njun Njxj Ntxt TNtX TVyX} \\
 $1$ & \texttt{Vyxt VzYU Vzun Vzxj YuNJ YuVZ jYun juNn jxJn txTn} \\
 $1$ & \texttt{uNNJ uNVZ uNjy uVzy xJNJ xJVZ xTNT xTVY yxTv zYuv} \\
 $1$ & \texttt{zuNv zxJv} \\

\end{tabular}
\end{center}

There are 179 entries, with distinct three-letter inputs.  In each
entry the number of capital input letters is one more than the
number of capital output letters.  This proves the weight assertion.

Here are explicit finite recurrences for checking the remaining
identities.  They also specify how every zero entry participates in
the calculation.  Let $\mathcal B$ be the twenty-letter basis and
$\mathcal B_{\mathfrak r}=\mathcal B\setminus\{e,f\}$.  Write
\[
 ab=\sum_v\mu_{ab}^{v}v,
 \qquad p(a,b,c)=\sum_v p_{abc}^{v}v.
\]
The lower-case coefficients $\mu$ are exactly
\eqref{alg:C-products} and the idempotent actions.  The remaining
coefficients are determined, for lower-case $a,b,c$, by
\begin{equation}
 \mu_{a,b^*}^{c^*}=\mu_{ca}^{b},\qquad
 \mu_{b^*,a}^{c^*}=\mu_{ac}^{b},\qquad
 \mu_{a^*,b^*}^{v}=0.
\label{coc:transpose-recurrence}
\end{equation}
Unspecified coefficients are zero.  In particular every coefficient
belongs to $\mathbb F_2[q]$.

For each $a,b,c,d\in\mathcal B_{\mathfrak r}$ and $v\in\mathcal B$,
substitution of the table and \eqref{coc:transpose-recurrence} gives
\begin{equation}
\begin{split}
0={}&\sum_{x\in\mathcal B}\bigl(
 \mu_{ax}^{v}p_{bcd}^{x}
 +\mu_{ab}^{x}p_{xcd}^{v}
 +\mu_{bc}^{x}p_{axd}^{v}\\
 &\hspace{35mm}
 +\mu_{cd}^{x}p_{abx}^{v}
 +p_{abc}^{x}\mu_{xd}^{v}\bigr).
\end{split}
\label{coc:closure-recurrence}
\end{equation}
This is the coefficient of $v$ in the normalized Hochschild
differential.  One may check it on all $18^4=104\,976$ radical
four-words; precisely $15\,250$ have composable endpoints, and the
other equations vanish by the corners.  Thus the finite calculation
covers the whole relative bar term, not a selected collection of
words.

For completeness, the boundary identity has an equally short
coefficient check.  Set $\epsilon_a=0$ on lower-case letters and
$\epsilon_a=1$ on capital letters.  For every basis letter put
\[
 \gamma_a=
 \begin{cases}
 q,&a\in\{u,v,t,E,X,Y,Z,U,V\},\\
 0,&\text{otherwise}.
 \end{cases}
\]
For all $a,b\in\mathcal B_{\mathfrak r}$ and $v\in\mathcal B$, the
same table gives
\begin{align}
 \sum_{\substack{w\in\mathcal B_{\mathfrak r}\\\epsilon_w=0}}
       \sum_x p_{abw}^{x}\mu_{x,w^*}^{v}
 &=\epsilon_b\gamma_a\mu_{ab}^{v},
 \label{coc:boundary-constant}\\
 \sum_{\substack{w\in\mathcal B_{\mathfrak r}\\\epsilon_w=1}}
       \sum_x p_{abw}^{x}\mu_{x,w^*}^{v}
 &=\bigl(\gamma_b+\gamma_v+(1-\epsilon_b)\gamma_a\bigr)
       \mu_{ab}^{v}.
 \label{coc:boundary-linear}
\end{align}
There are $18^2=324$ input pairs.  These two equations are the
constant and linear coefficients in $H$ of
\eqref{coc:boundary}: its left side gives $H^{\epsilon_w}$, while
the three terms on the right give respectively
$H\gamma_b\mu_{ab}^{v}$, $H\gamma_v\mu_{ab}^{v}$, and
$H^{1-\epsilon_b}\gamma_a\mu_{ab}^{v}$.  This verifies the identity
for an indeterminate $H$, hence for either $H_1$ or $H_2$.

All these substitutions are polynomial identities in characteristic
two.  They require no specialization of $q$ or $H$.  Finally the
entries $p(t,x,J)=0$ and $p(t,y,J)=q^3f$ give
\[
 q^2p(t,x,J)+p(t,y,J)=q^3f,
\]
which is the nonzero value on the explicit cycle
\eqref{coc:cycle}.
